% SPDX-FileCopyrightText: 2026 Will Cook
% SPDX-License-Identifier: CC-BY-4.0
%
% Standalone Erdős Problem Note for #257.  Context is deliberately repeated
% from the gateway so that this file remains independently readable.
%
% Ordering rule for this note: settled supports first, then the proved
% limitation, then the open frontier.  An earlier draft opened with a research
% programme and a list of analytic targets; a reader who stops after one page
% then sees only a plan, which is exactly the reading the note must not
% invite.
\documentclass[11pt]{article}
% SPDX-FileCopyrightText: 2026 Will Cook
% SPDX-License-Identifier: CC-BY-4.0
%
% Shared preamble for the Erdős Problem Notes: the short, standalone,
% problem-owned manuscripts covering the expansion library `ErdosProblems`.
%
% One note owns one Erdős problem.  Each note repeats whatever context its own
% reader needs, so that a reader who arrives at a single problem never has to
% assemble it from the gateway paper.  Deliberate repetition across notes is the
% design, not an oversight.
%
% Source links resolve against one pinned formal-source commit, declared once
% below.  `scripts/check_problem_note_sources.py` verifies that every authored
% (file, line, declaration) triple names that exact declaration in the pinned
% snapshot, so a link cannot silently rot as later work moves lines.

\usepackage{paper-house-style}
\usepackage{array}
\usepackage{booktabs}
\usepackage{enumitem}
\setlist{topsep=0.35em,itemsep=0.2em}
\usepackage[colorlinks=true,linkcolor=linkinternal,urlcolor=linkexternal,
  citecolor=linkinternal]{hyperref}
\hypersetup{bookmarksnumbered=true,bookmarksopen=true,bookmarksopenlevel=1,
  pdfauthor={Will Cook},pdflang={en-GB}}

% ---------------------------------------------------------------------------
% Pinned formal source.  \commit names the checkpoint every link resolves
% against; the checker reads that snapshot rather than the working tree, so the
% printed coordinates stay correct however later waves move declarations.
% ---------------------------------------------------------------------------
\newcommand{\commit}{e5fd26a6d1b1a346430205eab5385b4c9565d004}
\newcommand{\commitshort}{e5fd26a6d1b1}
\newcommand{\repobase}{https://github.com/wcook04/plectis-lean-erdos249-257/blob/\commit}
\newcommand{\sourceurl}{https://github.com/wcook04/plectis-lean-erdos249-257/tree/\commit}
\newcommand{\PX}{ErdosProblems}

% Declaration names stay inside link targets.  The printed label is either a
% spaced, lower-case reading of the declaration name (\lref) or a short authored
% phrase (\lword); neither prints a module path or a raw Lean identifier.
\ExplSyntaxOn
\cs_new_protected:Npn \noteleanlabel #1
  {
    \tl_set:Nx \l_tmpa_tl { \tl_to_str:n {#1} }
    \regex_replace_all:nnN { _+ } { \c{space} } \l_tmpa_tl
    \regex_replace_all:nnN
      { ([a-z0-9])([A-Z]) }
      { \1\c{space}\2 }
      \l_tmpa_tl
    \tl_set:Nx \l_tmpa_tl { \tl_lower_case:n { \tl_use:N \l_tmpa_tl } }
    \tl_use:N \l_tmpa_tl
  }
\ExplSyntaxOff
\newcommand{\leanlabel}[1]{\noteleanlabel{#1}}
\newcommand{\lref}[3]{\href{\repobase/\PX/#1\#L#2}{\leanlabel{#3}}}
\newcommand{\lword}[4]{\href{\repobase/\PX/#1\#L#2}{#4}}
\newcommand{\lloc}[2]{\href{\repobase/\PX/#1\#L#2}{Lean source}}

% Links into a sibling library, named repository-relative.  The expansion
% library is implicit in \lref/\lword, because that is where problem-owned
% work lands.  Several problems have their headline theorem in the older
% reviewed #249/#257 corpus, and a note that could not name that library
% would have to paraphrase a checked statement instead of linking it.
\newcommand{\mref}[3]{\href{\repobase/#1\#L#2}{\leanlabel{#3}}}
\newcommand{\mword}[4]{\href{\repobase/#1\#L#2}{#4}}
\newcommand{\mloc}[2]{\href{\repobase/#1\#L#2}{Lean source}}
\emergencystretch=2em

\newtheorem{theorem}{Theorem}[section]
\newtheorem{proposition}[theorem]{Proposition}
\newtheorem{lemma}[theorem]{Lemma}
\newtheorem{corollary}[theorem]{Corollary}
\theoremstyle{definition}
\newtheorem{definition}[theorem]{Definition}
\newtheorem{example}[theorem]{Example}
\newtheorem{problem}[theorem]{Problem}
\theoremstyle{remark}
\newtheorem*{remark}{Remark}

\newcommand{\N}{\mathbb{N}}
\newcommand{\Npos}{\mathbb{N}_{>0}}
\newcommand{\Z}{\mathbb{Z}}
\newcommand{\Q}{\mathbb{Q}}
\newcommand{\R}{\mathbb{R}}
\newcommand{\lcm}{\operatorname{lcm}}
\newcommand{\ph}{\varphi}

% The series line printed under every note's title block.
\newcommand{\seriesline}{Erd\H{o}s Problem Notes}

% Production provenance.  The wording is fixed by the corpus provenance
% contract and reproduced verbatim in every manuscript in this repository, so
% the papers cannot drift into different accounts of how they were made.
\newcommand{\noteprovenance}{\provenancenote{\emph{Authorship and AI use.} The
prose, formal proofs, and software in this version were drafted and revised by
large-language-model agents under Will Cook's direction.  Cook set the
objectives and acceptance criteria, reviewed and authorised the manuscript, and
is responsible for its contents.  The agents are production tools, not authors.
See \emph{Statements and declarations}.}}

% The status paragraph every note carries, in its own words, before any result.
% Kernel checking establishes the exact Lean proposition; it does not establish
% that the proposition is the interesting one, that it is new, or that the
% Erdős problem is closed.
\newcommand{\statusboundary}{%
  \emph{Status.}  The problem treated here is open, and this note does not
  close it.  Every statement below marked as checked is a proposition that the
  pinned Lean kernel accepts from the sources this note links to, with no
  \texttt{sorry}, no added axiom, and no unchecked evaluation.  That is a claim
  about the formal statement, not about its mathematical interest, its
  novelty, or the original problem.  The unresolved obligations are named
  exactly, in their own section, and none of the finite computations,
  reductions, or no-go results here removes one of them.}

\hypersetup{
  pdftitle={Settled supports and exact finite periods: a machine-checked map for Erdős Problem 257},
  pdfsubject={Erdős Problem Notes: Problem 257},
  pdfkeywords={irrationality, Mersenne numbers, achievement sets, divisor incidence, squarefree, Lean 4},
}

\newcommand{\wgt}[1]{w_{#1}}
\newcommand{\Ach}{\mathcal A}
\newcommand{\scA}{\operatorname{sc}}

\runninghead{Erd\H{o}s \#257: supports and finite periods}
\title{Settled Supports and Exact Finite Periods}
\subtitle{A machine-checked map for Erd\H{o}s Problem~\#257, with a
squarefree-support method boundary}
\author{Will Cook}
\date{July 2026}

\begin{document}
\maketitle
\noteprovenance

\begin{abstract}
\noindent
Erd\H{o}s~\#257 asks whether $X_A=\sum_{a\in A}(2^a-1)^{-1}$ is irrational for
every infinite $A\subseteq\Npos$.  It remains open.  This note separates the
content already proved from the universal endpoint.

For every finite nonempty exponent set $F$ and every integer base $b\ge2$, the
multiplicative order of $b$ modulo the reduced denominator of
$\sum_{n\in F}(b^n-1)^{-1}$ is exactly $\operatorname{lcm}(F)$: reduction never
collapses the period.  For infinite supports, seven structured families are
irrational in every base $b\ge2$, and periodic integer weights satisfy an
irrational-or-base-terminating dichotomy.  These are unconditional theorems,
not reformulations of universal \#257; their prior-art boundaries are stated
where they occur.

The note also records the achievement-set geometry and exact open dichotomy at
$1/2$, then proves a method limitation on the squarefree support.  Its divisor
incidence is $2^{\omega(n)}-1$, hence odd for every $n\ge2$, so the two
certificate engines requiring first-block divisibility have no instance
there.  This does not decide the squarefree value; it names exactly which
methods fail and which weaker engine survives.
\end{abstract}

\statusboundary

\section{The problem}\label{sec:problem}

\begin{problem}[Erd\H{o}s \#257]\label{res:problem}
For every infinite $A\subseteq\Npos$, is
$X_A=\sum_{a\in A}\dfrac1{2^a-1}$ irrational?
\end{problem}

\noindent Numbering and status follow
\href{https://www.erdosproblems.com/257}{Bloom's catalogue}.  Write
$\wgt n=(2^n-1)^{-1}$.  Interchanging sums gives the coordinate this note
works in: with the \emph{divisor incidence}
$\scA_A(n)=\#\{d\mid n: d\in A\}$,
\[
  X_A=\sum_{a\in A}\frac1{2^a-1}=\sum_{n\ge1}\frac{\scA_A(n)}{2^n}.
  \tag{1}\label{eq:incidence}
\]
So \#257 is the question of which $0/1$ weights make a Lambert series
irrational, and every argument below reads the coefficient sequence
$\scA_A$ rather than $A$ itself.

\section{Exact period of every finite partial sum}\label{sec:period}

For a finite nonempty $F\subseteq\Npos$ and an integer $b\ge2$, write
\[
  x_F(b)=\sum_{n\in F}\frac1{b^n-1}=\frac{N_F}{D_F}
  \qquad (\gcd(N_F,D_F)=1,\ D_F>0).
\]

\begin{theorem}[finite-period noncollapse]\label{res:period}
The reduced denominator is coprime to $b$, and
\[
  \operatorname{ord}_{D_F}(b)=\operatorname{lcm}\{n:n\in F\}.
\]
Thus cancellation in the numerator never lowers the base-$b$ period below
the least common multiple of the selected exponents.
\end{theorem}

\noindent The general theorem is
\mword{Erdos249257/CertificateKernel.lean}{5091}{finite_period_noncollapse}{checked};
the reduced-denominator and coprimality forms are
\mword{Erdos249257/CertificateKernel.lean}{5246}{finite_period_noncollapse_rat_den}{checked}
and
\mword{Erdos249257/CertificateKernel.lean}{5061}{coprime_base_reducedDenominator}{checked}.
This is an unconditional statement about every finite support, not a bounded
table of examples and not an implication with an open hypothesis.  It does not
settle any infinite support.  Denominator-period control is part of the
classical engine behind Erd\H{o}s's 1948 argument; whether this exact sharp
noncollapse statement appears in the literature has not been assessed, so no
novelty claim is made here.

\section{The map of settled supports}\label{sec:map}

\begin{table}[h]
\centering
\small
\begin{tabular}{@{}
  >{\raggedright\arraybackslash}p{0.31\linewidth}
  >{\raggedright\arraybackslash}p{0.15\linewidth}
  >{\raggedright\arraybackslash}p{0.46\linewidth}@{}}
\toprule
Support $A$ & Bases & Authority \\
\midrule
all of $\Npos$ & every $b\ge2$ & Erd\H{o}s 1948;
  \mword{Erdos249257/CertificateKernel.lean}{8000}{irrational_erdosSum_full_support}{checked here} \\
primes & $b=2$ (and $b\ge2$) & Tao--Ter\"av\"ainen~\cite{taoteravainen2025};
  not formalised here \\
prime powers & $b=2$ & \emph{ibid.}, stated as a remark; not formalised here \\
factorials $\{n!\}$ & every $b\ge2$ &
  \mword{Erdos249257/CertificateKernel.lean}{5754}{erdos257_family_factorial_instance}{checked here} \\
powers of two $\{2^k\}$ & every $b\ge2$ &
  \mword{Erdos249257/CertificateKernel.lean}{5762}{erdos257_family_two_pow_instance}{checked here} \\
multiples of a fixed $d$ & every $b\ge2$ &
  \mword{Erdos249257/CertificateKernel.lean}{8775}{irrational_erdosSupportSeries_multiples}{checked here} \\
a residue class & every $b\ge2$ &
  \mword{Erdos249257/CertificateKernel.lean}{11344}{irrational_erdosSupportSeries_residueClass}{checked here} \\
the odd numbers & every $b\ge2$ &
  \mword{Erdos249257/CertificateKernel.lean}{11358}{irrational_erdosSupportSeries_odd}{checked here} \\
eventually periodic & every $b\ge2$ &
  \mword{Erdos249257/CertificateKernel.lean}{11276}{irrational_erdosSupportSeries_eventuallyPeriodic}{checked here} \\
pairwise coprime, $\sum_{a\in A}a^{-1}<\infty$ & every $b\ge2$ &
  \mword{Erdos249257/CertificateKernel.lean}{10448}{irrational_erdosSupportSeries_pairwise_coprime}{checked here} \\
\midrule
squarefree numbers & --- & \textbf{Open};
  outside two divisibility-first engines here (Section~\ref{sec:squarefree}) \\
arbitrary infinite $A$ & --- & \textbf{Open} (Problem~\ref{res:problem}) \\
\bottomrule
\end{tabular}
\end{table}

\noindent Three things about this table are worth stating plainly.

First, the families checked here are proved for \emph{every} integer base
$b\ge2$, not only base $2$; the general statement is
\mword{Erdos249257/CertificateKernel.lean}{8717}{irrational_erdosSupportSeries_univ}{the base-$b$ full-support theorem}
and the family results specialise it.  Prior art is acknowledged rather than
absorbed: for eventually periodic \emph{rational} coefficient sequences Luca
and Tachiya prove a broader theorem by different methods~\cite{lucatachiya},
and no priority is claimed for the formalised versions.

Second, the two method families are \emph{complementary, not nested}.  The
analytic method of~\cite{taoteravainen2025} reaches the primes and prime
powers, which are neither eventually periodic nor pairwise coprime with
summable reciprocals, so they are outside every row proved here.  The rows
proved here --- factorials, powers of two, residue classes --- are not
``sufficiently similar to the primes'' in the sense that method needs.
Neither list contains the other, and their union does not exhaust the
infinite supports.

Third, none of this approaches the universal statement.  A list of settled
families, however long, is a list; Problem~\ref{res:problem} quantifies over
all infinite $A$, and the gap between the two is not narrowed by adding rows.

\begin{theorem}[periodic signed-weight dichotomy]\label{res:signed}
Let $u:\Npos\to\mathbb Z$ be periodic and $b\ge2$.  Then
\[
  \sum_{n\ge1}\frac{u(n)}{b^n-1}
\]
is irrational, or some power of $b$ times this value is an integer.
\end{theorem}

\noindent This is
\mword{Erdos249257/CertificateKernel.lean}{13847}{irrational_or_bpow_mul_eq_intCast_intWeightedErdosSeries_periodic}{checked};
the source also proves the full-block certificate supply used by the theorem.
For nonnegative support indicators the terminating alternative is excluded in
the settled rows above.  For genuinely mixed signs it is not excluded in
general, so Theorem~\ref{res:signed} is a dichotomy rather than a blanket
irrationality theorem.  Prior art for this mixed-sign periodic formulation has
not been assessed.

\section{A support outside two certificate engines}\label{sec:squarefree}

Let $A_{\mathrm{sf}}=\{d\ge2: d\text{ squarefree}\}$.  This is a natural
support of density $6/\pi^2$ --- far denser than the primes, and not periodic.

\begin{theorem}[squarefree divisor incidence]\label{res:sfcount}
For every $n\ge1$ the number of squarefree divisors of $n$ is $2^{\omega(n)}$,
where $\omega(n)$ is the number of distinct prime factors, and hence
\[
  \scA_{A_{\mathrm{sf}}}(n)=2^{\omega(n)}-1 .
\]
In particular $\scA_{A_{\mathrm{sf}}}(n)$ is \emph{odd} for every $n\ge2$.
\end{theorem}

\noindent The count is
\lword{Erdos257/SquarefreeSupportIncidence.lean}{53}{card_squarefreeDivisors}{checked},
the incidence formula is
\lword{Erdos257/SquarefreeSupportIncidence.lean}{70}{squarefreeIncidence_eq}{checked},
and the parity conclusion is
\lword{Erdos257/SquarefreeSupportIncidence.lean}{100}{odd_squarefreeIncidence}{checked}.
The proof is the bijection between squarefree divisors of $n$ and subsets of
its prime factors
(\lword{Erdos257/SquarefreeSupportIncidence.lean}{30}{squarefreeDivisors_eq_image}{Lean}),
so the count is $2^{\omega(n)}$; removing $d=1$ leaves an odd number whenever
$\omega(n)\ge1$.

\begin{corollary}[two block-certificate engines cannot see this support]\label{res:blind}
For the squarefree support and every even base $b\ge2$, the digitwise and
carry-aware block-certificate schemas have no instance already at precision
$q=b^2$.  Their first-block divisibility condition would require $b$ to divide
the odd coefficient $\scA_{A_{\mathrm{sf}}}(N+1)$.
\end{corollary}

\noindent The current source contains direct Lean declarations for both
exclusions.  They postdate this problem-note series' pinned hyperlink snapshot,
so they are not hyperlinked here.  The argument does \emph{not} reach the
divisibility-free low-carry engine, odd bases, the support with $1$ adjoined,
or methods based on lcm gaps or near-integer estimates.

\noindent We are deliberate about what this does and does not say.

\emph{It does not say} that $X_{A_{\mathrm{sf}}}=\sum_{d\ \mathrm{squarefree},\,d\ge2}(2^d-1)^{-1}$
is rational, or that it is hard, or that it is irrational.  That value's status
is open and this note does not move it.

\emph{It does say} that two fully specified proof schemas provably cannot
certify one natural support at even bases, for a reason internal to the support
rather than to the search.  The obstruction is a parity accident rather than
a scaling one: the shifted coefficient $2^{\omega(n)}$ still satisfies
$2^{\omega(n)}\le n$
(\lword{Erdos257/SquarefreeSupportIncidence.lean}{117}{two_pow_card_primeFactors_le}{checked}),
so it remains inside the coefficient range those arguments require.  Repairing
the parity is therefore not obviously impossible --- but a repaired first block
does not discharge the remaining tail and head obligations, and we make no
claim that it does.

\section{Achievement-set geometry}\label{sec:geometry}

The set of all values obtainable from the base-$2$ Mersenne weights is the
achievement set
$\Ach=\{\sum_{n\ge1}\varepsilon_n\wgt n:\varepsilon_n\in\{0,1\}\}$.

\begin{theorem}[geometry of $\Ach$]\label{res:geometry}
$\Ach$ is compact, perfect, totally disconnected and nowhere dense, and has
Lebesgue measure~$1$.
\end{theorem}

\noindent Checked as
\mword{Erdos249257/GreedyAchievementSet.lean}{1620}{perfect_mersenneAchievementSet}{perfect},
\mword{Erdos249257/GreedyAchievementSet.lean}{1636}{isTotallyDisconnected_mersenneAchievementSet}{totally disconnected},
\mword{Erdos249257/GreedyAchievementSet.lean}{1645}{isNowhereDense_mersenneAchievementSet}{nowhere dense},
and
\mword{Erdos249257/GreedyAchievementSet.lean}{996}{volume_mersenneAchievementSet}{of measure one}.
The strict-tail Cantor geometry is recorded by Kova\v{c} and Tao; no novelty is
claimed for the formalised refinements.  Membership is equivalent to survival
of the greedy expansion at every level
(\mword{Erdos249257/GreedyAchievementSet.lean}{1445}{mem_mersenneAchievementSet_iff_greedy_survival}{Lean}),
and exact rational death certificates prove non-membership: $3/4\notin\Ach$
(\mword{Erdos249257/GreedyAchievementSet.lean}{1748}{three_fourths_not_mem_mersenneAchievementSet}{Lean}).

\subsection{Rigidity is hereditary}\label{sec:hereditary}

Problem~\ref{res:problem} quantifies over \emph{every} infinite support, so it
matters whether the base-$2$ coding survives deleting weights.  It does, and
the argument is short.  For any set $J$ of future offsets, write
$T_J(n)=\sum_{k\in J}\wgt{n+k+1}$ for the tail restricted to $J$.

\begin{theorem}[hereditary strict superincreasingness]\label{res:hereditary}
For every $J\subseteq\N$ and every $n\ge1$, $T_J(n)<\wgt n$.  Consequently the
digit map is injective on strings supported in $J$: every Mersenne subseries
retains unique binary coding.
\end{theorem}

\noindent Checked as
\lword{Erdos257/MersenneSubseriesRigidity.lean}{27}{selectedMersenneTail_lt_weight}{the hereditary inequality}
and
\lword{Erdos257/MersenneSubseriesRigidity.lean}{51}{supportedMersenneDigitValue_injective}{injectivity on a chosen support},
with summability of the restricted tail at
\lword{Erdos257/MersenneSubseriesRigidity.lean}{21}{summable_selectedMersenneTail}{Lean}.
The coding is stated on the subtype of
\lword{Erdos257/MersenneSubseriesRigidity.lean}{42}{SupportedMersenneDigits}{digit strings vanishing off $J$},
evaluated by the
\lword{Erdos257/MersenneSubseriesRigidity.lean}{46}{supportedMersenneDigitValue}{restricted digit map},
so the statement is about the subseries itself rather than a projection of the
full one.
The content is that strict superincreasingness is inherited by \emph{arbitrary}
subfamilies, not merely by tails: deleting weights can only shrink $T_J(n)$,
which already sits below $\wgt n$ at the full support.  So no choice of
infinite $A$ escapes the rigidity used throughout this note, and a
counterexample to \#257 cannot be sought by weakening the coding.  This
constrains where a counterexample could live; it does not produce one, and it
decides no value.

\section{The half-value dichotomy}\label{sec:half}

A rational point of $\Ach$ attained by an infinite support would refute
Problem~\ref{res:problem}.  The distinguished candidate is $1/2$.

\begin{theorem}[exact dichotomy at $1/2$]\label{res:half}
$1/2\in\Ach$ if and only if the canonical greedy expansion of $1/2$ omits
infinitely many exponents.  Moreover no finite support has value $1/2$, so
membership would automatically produce an \emph{infinite} support of rational
value, and hence a counterexample to universal \#257.
\end{theorem}

\noindent The equivalence is
\mword{Erdos249257/GreedyAchievementSet.lean}{2514}{half_mem_mersenneAchievementSet_iff_greedySkippedSupport_infinite}{checked},
and the finite-support exclusion is
\mword{Erdos249257/HalfCarryReachability.lean}{589}{finite_boolSupport_ne_half}{checked}.

This is an equivalence, and we mark it as such rather than as progress: both
sides are open, and proving them equivalent does not decide either.  It is
included because it fixes what a counterexample would have to look like, and
because the negative branch is genuinely one-sided --- non-membership is
witnessed by a finite fatal gap and is therefore semi-decidable, whereas
membership is an infinite conjunction that no finite computation certifies.
Computing the greedy orbit further can only ever raise a lower bound on where
death could occur; it cannot establish survival.

\section{What we need from a mathematician}\label{sec:ask}

Three points in this development stop for mathematical rather than
engineering reasons.  Each is set out below as a specification: the statement
wanted, what is already checked around it, which kind of answer would help
most, and exactly what would follow from one.

Two remarks apply to all three.  Each item is self-contained; answering one
needs no engagement with the other items, with the Lean sources, or with the
rest of this note, and an answer in ordinary mathematics is what is being
asked for.  And a negative answer is worth as much as a positive one: a proof
that a named route cannot work, or a reference showing it is a known dead
end, would be acted on by withdrawing the route from this note rather than by
defending it.

\subsection{A method reaching a support outside both lists}\label{sec:ask:univ}

\emph{Statement wanted.}  An infinite $A\subseteq\Npos$, with a proof that
$X_A$ is irrational, such that $A$ is not eventually periodic, is not
pairwise coprime with $\sum_{a\in A}a^{-1}<\infty$, and is neither the primes
nor the prime powers; better, a theorem covering an infinite class of such
$A$.  Every family checked here is produced, directly or after a base change,
by one engine, whose only hypothesis is that for each precision $q$ the
incidence sequence $\scA_A$ admits a block certificate.  That hypothesis is
the premise of
\mword{Erdos249257/CertificateKernel.lean}{8703}{irrational_erdosSupportSeries_of_weighted_coeff_certificates}{irrational erdos support series of weighted coeff certificates}
and of its carry-aware form
\mword{Erdos249257/CertificateKernel.lean}{9053}{irrational_erdosSupportSeries_of_weighted_coeff_carry_certificates}{checked}.
A method that does not factor through certificates at all is equally
admissible, and would be more informative.

\emph{Already checked.}  The reduction from irrationality to certificate
supply is proved for every base $b\ge2$; the supply itself, for a general
$A$, is what is missing, and the theorem statements say so.  The supports for
which it has been produced are the rows of Section~\ref{sec:map}, and that
section also records that they neither contain nor are contained in the
primes and prime powers of~\cite{taoteravainen2025}.  Two constraints on any
answer are already fixed: every finite partial sum has base-$b$ period exactly
the least common multiple of its exponents (Theorem~\ref{res:period}), and
strict superincreasingness is inherited by every subfamily
(Theorem~\ref{res:hereditary}), so no counterexample can be obtained by
weakening the binary coding.

\emph{What would help most.}  A reason no single method can reach every
infinite $A$: for instance an infinite $A$ whose series is irrational but
which provably admits no certificate of the form above.  That would say the
engine has the wrong shape rather than too little fuel, and it is the answer
we are least able to produce ourselves.  Next best, a construction reaching
one support outside both lists.  A reference doing either is worth as much as
doing it.

\emph{If it were answered.}  Neither answer closes Problem~\ref{res:problem}.
A construction adds one row to Section~\ref{sec:map} and settles only the
class it reaches; as that section says, a list of families does not approach
a universal statement.  A no-single-method result settles no case of the
problem at all, and would still be the more useful of the two, because it
would say which programmes to abandon.

\subsection{The squarefree support, without first-block divisibility}
\label{sec:ask:sf}

\emph{Statement wanted.}  Decide
$X_{A_{\mathrm{sf}}}=\sum_{n\ge1}\bigl(2^{\omega(n)}-1\bigr)2^{-n}$ by an
argument that never requires a first block of coefficients to be divisible by
the base.  Writing $f(n)=2^{\omega(n)}-1$, the engine that survives
Corollary~\ref{res:blind} asks, for every $q\ge1$, for $N$, $K\le L$ and $C$
with
\[
  \sum_{r=K+1}^{L}f(N+r)\,2^{L-r}\le C,
  \qquad
  V=\Bigl(\sum_{r=1}^{K}f(N+r)\,2^{K-r}\Bigr)\bmod 2^{K},
\]
some nonzero coefficient beyond $N+L$, and
\[
  q\bigl(V\cdot 2^{L-K}+C+N+L+2\bigr)<2^{L}.
\]
Nothing is asked of the first block except that its carry residue $V$ be
small enough to be absorbed.
That premise is the hypothesis of
\mword{Erdos249257/CertificateKernel.lean}{14707}{irrational_coeff_series_of_weighted_coeff_low_carry_block_certificates}{irrational coeff series of weighted coeff low carry block certificates}.

\emph{Already checked.}  The incidence identity $f(n)=2^{\omega(n)}-1$ and
its oddness for every $n\ge2$ (Theorem~\ref{res:sfcount}), and the resulting
exclusion of the two divisibility-first engines at every even base
(Corollary~\ref{res:blind}).  The growth envelope this engine needs,
$f(n)\le n$, follows from the bound $2^{\omega(n)}\le n$ recorded in
Section~\ref{sec:squarefree}, so the coefficient sequence is admissible and
only the certificate supply is missing.  The value itself is open in both
directions; nothing here bears on which way it goes.

\emph{What would help most.}  A reference.  The function $2^{\omega}$ is
standard, and $\sum_{n\ge1}2^{\omega(n)}2^{-n}$ differs from
$X_{A_{\mathrm{sf}}}$ by $\sum_{n\ge1}2^{-n}=1$, so either determines the
other; if that Lambert-type series has already been treated, the pointer
retires this item outright.  Failing that, a construction: an explicit supply
of low-carry certificates for $f$, or any analytic route to the value.

\emph{If it were answered.}  Irrationality of $X_{A_{\mathrm{sf}}}$ closes
the open row in Section~\ref{sec:map} and item~2 of Section~\ref{sec:open}.
It does not touch Problem~\ref{res:problem}, which quantifies over all
infinite supports.  Rationality would answer Problem~\ref{res:problem} in the
negative, the squarefree numbers being infinite; we have no evidence for
that, and record it only because the item is two-sided.

\subsection{The two surviving cells at the half value}\label{sec:ask:half}

By Theorem~\ref{res:half}, $1/2\in\Ach$ exactly when the greedy expansion of
$1/2$ omits infinitely many exponents.  Suppose it does not.  In the
coordinates this development works in, the expansion is then eventually right
and has a last non-right transition at some rank $D\ge13$
(\mword{Erdos249257/HalfCylinderLastProducerContradiction.lean}{178}{exists_last_false_terminal_of_eventuallyRight}{checked}),
of one of two kinds.  The first kind is excluded unconditionally
(\mword{Erdos249257/HalfCylinderLastProducerContradiction.lean}{250}{upperProducer_not_last}{checked}).
The second carries an integer $c(D)$, its producer carry, and the surviving
obligation is that the complete future divisor-incidence tail from row $2D+2$
lies strictly below $c(D)$.  That tail is nonnegative, so a cell with
$c(D)<0$ cannot be settled by the comparison and has to be excluded outright.
Cell $-3$ has been excluded
(\mword{Erdos249257/HalfCylinderLastProducerContradiction.lean}{315}{middleProducer_neg_three_not_last}{checked});
cells $-2$ and $-1$ have not.

\emph{Statement wanted.}  The tail comparison at every last middle transition
with $D\ge13$ and $c(D)\ne-3$.  That statement is
\mword{Erdos249257/HalfCylinderLastProducerContradiction.lean}{54}{SeamMiddleProducerTailEscapeExceptNegThree}{seam middle producer tail escape except neg three},
and it contains three separable pieces of work.

\begin{enumerate}
\item Cell $-2$ does not occur.  What is checked is a conditional bound: when
  the lazy support --- the finite word below $D$ together with every rank
  above $D$ --- has value below $1/2$, a cell $-2$ forces its
  divisor-incidence tail at row $2D+2$ below $2$
  (\mword{Erdos249257/HalfCylinderFinalMiddleTailSocket.lean}{56}{binaryCoeffTail_union_Ioi_lt_two_of_producerCarry_eq_neg_two}{Lean}).
  It is therefore enough to show that same tail is at least $2$, which is
  done whenever $D\not\equiv2\pmod3$
  (\mword{Erdos249257/HalfFinalMiddlePhaseSieve.lean}{50}{two_lt_binaryCoeffTail_union_Ioi_of_forced_two_three_six}{Lean}).
  A further sieve through rank $26$ leaves $412$ joint residue classes modulo
  $2730$
  (\mword{Erdos249257/HalfFinalMiddlePhaseSieve.lean}{2050}{finalMiddleTwentySixPhaseSurvivors_card}{Lean}),
  and those are not excluded.
\item Cell $-1$ does not occur.  The same conditional bound holds with
  threshold $3$
  (\mword{Erdos249257/HalfCylinderFinalMiddleTailSocket.lean}{74}{binaryCoeffTail_union_Ioi_lt_three_of_producerCarry_eq_neg_one}{Lean}),
  so a tail of at least $3$ would settle it.  No comparable sieve is recorded
  for this cell.
\item The comparison itself at every remaining cell: at every $c(D)\ge0$, and
  at any cell below $-3$ should one occur.  A cruder sufficient form asks only
  $c(D)>2\sqrt{2D+2}+4$
  (\mword{Erdos249257/HalfCylinderLastProducerContradiction.lean}{84}{SeamMiddleProducerSqrtEscape}{Lean}).
\end{enumerate}

\emph{What would help most.}  Item~3, proved or refuted.  Items~1 and~2 are
finite-looking and would shrink the obligation, but neither decides anything
on its own.  A demonstration that the comparison fails at infinitely many
middle transitions would show the route cannot work, and we would record that
here.

\emph{If it were answered.}  The three together give $1/2\in\Ach$
(\mword{Erdos249257/HalfCylinderLastProducerContradiction.lean}{387}{half_mem_mersenneAchievementSet_of_middleProducerTailEscapeExceptNegThree}{checked}),
and by Theorem~\ref{res:half} the support realising it is infinite, so
Problem~\ref{res:problem} would be answered in the negative.  We attach no
likelihood to that outcome: the hypothesis is open, and its failure would
prove nothing about Problem~\ref{res:problem} in either direction.  Taken
singly the three settle nothing: items~1 and~2 close one branch each and
decide no value, and item~3 leaves both cells standing.  A proof that
$1/2\notin\Ach$ would retire this item as well, and by Theorem~\ref{res:half}
it would take the form of a finite fatal gap in the greedy orbit.

\section{What remains open}\label{sec:open}

\begin{enumerate}
\item Problem~\ref{res:problem} itself, for arbitrary infinite $A$.  The map in
  Section~\ref{sec:map} is a list of families and does not approach a universal
  statement.
\item Irrationality of $X_{A_{\mathrm{sf}}}$ for the squarefree support.
  Section~\ref{sec:squarefree} closes one method against it and supplies no
  replacement.
\item Membership of $1/2$ in $\Ach$.  Theorem~\ref{res:half} makes it exactly
  equivalent to infinitely many greedy skips; neither side is proved.
\item Which supports the analytic method of~\cite{taoteravainen2025} reaches.
  The authors name this themselves and do not pursue it; it is, in our view,
  the most valuable open question in the immediate neighbourhood of \#257, and
  it is not one this development is currently equipped to answer.
\end{enumerate}

\section*{Statements and declarations}

This manuscript is authored exposition, not proof authority.  The linked Lean
snapshot is authoritative only for its exact propositions; kernel checking
establishes that a proposition was proved, not that it is interesting, novel,
or sufficient.  Results attributed to Tao and Ter\"av\"ainen, to Luca and
Tachiya, and to Kova\v{c} and Tao are cited from the literature and are not
formalised here.  Erd\H{o}s Problem~\#257 remains open.

\begin{thebibliography}{9}
\bibitem{taoteravainen2025} T.~Tao and J.~Ter\"av\"ainen,
  \emph{Quantitative correlations and some problems on prime factors of
  consecutive integers}, arXiv:2512.01739 (submitted December 2025, revised
  April 2026).  Proves $\sum_{n\ge1}\omega(n)/2^n=\sum_{p}(2^p-1)^{-1}$
  irrational, settling the prime-support case of \#257; remarks that the
  method also gives the prime-power support and every integer base $b\ge2$.
\bibitem{pratt} K.~Pratt, \emph{The irrationality of a prime factor series
  under a prime tuples conjecture}, arXiv:2409.15185.
\bibitem{lucatachiya} F.~Luca and Y.~Tachiya, \emph{Irrationality of Lambert
  series associated with a periodic sequence}.
\bibitem{erdos1948} P.~Erd\H{o}s, \emph{On arithmetical properties of Lambert
  series}, J.~Indian Math.\ Soc.\ 12 (1948), 63--66.
\bibitem{erdosproblems} T.~Bloom, \emph{Erd\H{o}s Problems},
  \url{https://www.erdosproblems.com/257}.
\end{thebibliography}

\end{document}
